\documentclass[11pt]{article}

\usepackage[spanish]{babel}
\usepackage[margin=1in]{geometry}
\usepackage{amsmath, amssymb, amsfonts}
\usepackage{booktabs}
\usepackage{enumitem}
\usepackage{tcolorbox}
\usepackage{hyperref}
\usepackage{listings}
\usepackage{xcolor}
\usepackage{parskip}
\usepackage{pgfplots}

\usetikzlibrary{arrows.meta}
\pgfplotsset{compat=1.18}

\lstset{
  language=Python,
  basicstyle=\ttfamily\footnotesize,
  breaklines=true,
  columns=flexible,
  keepspaces=true,
  commentstyle=\color{gray}\itshape,
  keywordstyle=\color{blue!70!black}\bfseries,
  stringstyle=\color{green!40!black},
  frame=L,
  xleftmargin=2em,
  showstringspaces=false,
  tabsize=4
}

\title{\textbf{Laboratorio 9: Redes Bayesianas \& Redes de Markov} \vspace{1em}\\
  Inteligencia Artificial - CC3085}
\author{José Antonio Mérida Castejón}
\date{\today}

\begin{document}

\maketitle

\subsection*{Task 1 - Teoría}
\textit{Responda las siguientes preguntas de forma clara, concisa y argumentada. No basta con
definir: se espera que usted conecte el concepto con sus implicaciones prácticas o estratégicas.}

\begin{enumerate}

  \item \textit{En clase discutimos cómo la regla de Bayes convierte un razonamieno de diagnóstico en uno causal. Con ello considere
    que usted diseña un sistema de IA para detectar fraudes en tarjetas de crédito. Defina una variable ``Evidencia Observable'' (B)
    y una variable de ``Causa'' (A). Adicionalmente, explique por qué, desde el punto de vista del modelado de datos, suele ser más
    fácil estimar en producción $P(Evidencia | Causa)$ que $P(Causa | Evidencia)$}


  \item \textit{Basado en lo que aprendió en clase, compare las Redes de Markov con las Redes Bayesianas. En los siguientes escenarios,
    recuerde que debe indicar y justificar qué tipo de red (Markov o Bayesiana) utilizaría para cada escenario:}

    \begin{enumerate}
      \item\textit{Necesita modelar un sistema de segmentación de imágenes médicas donde el estado de un píxel (tejido sano) o tumor
        depende fuertemente del estado de sus píxeles vecinos de forma simétrica.}

        En este caso, logramos identificar claramente la necesidad de incluir probabilidades causales de nuestro modelo. Según el enunciado,
        la probabilidad del estado de un píxel se ve sesgada (o causada por) el estado de los píxeles anteriores. En este caso, una Red de
        Markov no sería suficiente para captar la información que necesitamos.

      \item\textit{Necesita modelar un sistema de diagnóstico de fallas en el motor de un vehículo, donde ciertas piezas al romperse
        provocan lecturas anómalas en los sensores.}

        En este caso, no tenemos probabilidades causales si no que probabilidades de diferentes estados del sistema dados diferentes
        factores. Este es un caso clásico dónde podemos aplicar Redes de Markov.
    \end{enumerate}

  En una Red de Markov con $n = 50$ variables booleanas, cada variable puede tomar valores en 
$\{0, 1\}$, por lo que el número total de combinaciones posibles es $2^{50}$
(más de un cuadrillón de asignaciones). La constante de normalización 
 recorre \textit{todas} las asignaciones 
posibles del espacio conjunto, lo que exige evaluar explícitamente cada una de esas $2^{50}$ 
combinaciones. Esto representa un desafío computacional prohibitivo en la industria por tres razones 
fundamentales: 

\begin{itemize}

  \item El crecimiento es \textbf{exponencial en $\mathbf n$}, de modo que agregar una sola variable 
adicional \textit{duplica} el costo de cómputo, haciendo inviable el escalamiento a redes reales con 
cientos o miles de variables; 


  \item En aplicaciones industriales como modelado de lenguaje, visión por 
computadora o sistemas de recomendación, los grafos tienen estructuras densas donde la factorización 
no reduce suficientemente el espacio de búsqueda para que el cálculo exacto sea tratable en tiempo 
real. 

\item Los recursos de memoria y tiempo de procesamiento requeridos superan con creces la 
capacidad del hardware disponible, incluso en clústeres de alto rendimiento. Para evadir este cálculo 
exacto, se emplean dos familias de algoritmos aproximados: los \textbf{métodos de muestreo de Monte 
Carlo mediante cadenas de Markov (MCMC)}, como el \textit{Gibbs Sampling}, que generan muestras 
de la distribución conjunta sin necesidad de conocer $Z$ al explorar el espacio de estados de forma 
estocástica.

\end{itemize}

\end{enumerate}

\end{document}
